%==============================================================================
%  ENGR 2405 -- Electrical Circuits I
%  Lab Report Template
%
%  Compile with pdflatex (run twice so the running header settles).
%  Referenced images : schematic.png  scope_shot.png  ltspice_plot.png
%==============================================================================
\documentclass[11pt]{article}

% ---- packages ---------------------------------------------------------------
\usepackage[utf8]{inputenc}
\usepackage[T1]{fontenc}
\usepackage[margin=1in]{geometry}
\usepackage{amsmath,amssymb}
\usepackage{booktabs}
\usepackage{array}
\usepackage{siunitx}
\usepackage{graphicx}
\usepackage{float}
\usepackage{enumitem}
\usepackage{fancyhdr}
\usepackage{caption}
\usepackage{xcolor}
\usepackage{listings}
\usepackage{circuitikz}
\usepackage{tikz}
\usetikzlibrary{arrows.meta,positioning,calc}
\usepackage[hidelinks]{hyperref}

% ---- siunitx v2 / v3 compatibility ------------------------------------------
% \qty is siunitx v3; if the local install is still v2, fall back to \SI.
\ifdefined\qty\else\newcommand{\qty}[2]{\SI{#1}{#2}}\fi
\ifdefined\unit\else\newcommand{\unit}[1]{\si{#1}}\fi

% ---- table + notation defaults ----------------------------------------------
\renewcommand{\arraystretch}{1.15}
\newcommand{\degre}{^{\circ}}

% ---- SciLab listing style ---------------------------------------------------
\definecolor{codekw}{rgb}{0.00,0.30,0.70}
\definecolor{codecom}{rgb}{0.20,0.55,0.20}
\definecolor{codestr}{rgb}{0.64,0.08,0.08}

\lstdefinelanguage{Scilab}{
    morekeywords={
        function,endfunction,if,then,else,elseif,end,for,while,do,
        select,case,break,continue,return,clc,clear,disp,mprintf,
        printf,zeros,ones,eye,linspace,inv,det,sum,length,size,abs,
        sqrt,real,imag,exp,sin,cos,tan,atan,plot,plot2d,xlabel,ylabel,
        title,legend,grid,poly,roots,find,pause,error,warning
    },
    sensitive=true,
    morecomment=[l]{//},
    morecomment=[s]{/*}{*/},
    morestring=[b]",
    morestring=[b]',
}
\lstdefinestyle{scilab}{
    language=Scilab,
    basicstyle=\ttfamily\footnotesize,
    keywordstyle=\color{codekw}\bfseries,
    commentstyle=\color{codecom}\itshape,
    stringstyle=\color{codestr},
    numbers=left,
    numberstyle=\tiny\color{gray},
    numbersep=8pt,
    frame=single,
    rulecolor=\color{gray!50},
    breaklines=true,
    showstringspaces=false,
    tabsize=4,
    captionpos=b,
    xleftmargin=14pt,
}

% ---- LTspice netlist style --------------------------------------------------
\lstdefinelanguage{LTspice}{
    morekeywords={
        .tran,.op,.dc,.ac,.step,.param,.model,.lib,.include,.inc,
        .meas,.measure,.save,.backanno,.end,.ends,.subckt,.global,
        .ic,.nodeset,.options,.option,.func,.temp,.four,.noise,.tf
    },
    sensitive=false,
    morecomment=[l]{*},
    morecomment=[l]{;},
    morestring=[b]",
}
\lstdefinestyle{netlist}{
    language=LTspice,
    basicstyle=\ttfamily\footnotesize,
    keywordstyle=\color{codekw}\bfseries,
    commentstyle=\color{codecom}\itshape,
    frame=single,
    rulecolor=\color{gray!50},
    breaklines=true,
    tabsize=4,
    captionpos=b,
    xleftmargin=14pt,
}

% ---- plain console / runtime printout style ---------------------------------
\lstdefinestyle{console}{
    basicstyle=\ttfamily\footnotesize,
    frame=single,
    rulecolor=\color{gray!50},
    breaklines=true,
    tabsize=4,
    captionpos=b,
    xleftmargin=14pt,
}

% ---- running header ---------------------------------------------------------
\pagestyle{fancy}
\fancyhf{}
\lhead{ENGR 2405 -- Lab <N>}
\rhead{Abdon Morales}
\cfoot{\thepage}
\renewcommand{\headrulewidth}{0.4pt}

%==============================================================================
\begin{document}

% ----------------------------- title block -----------------------------------
\begin{center}
    {\large\bfseries ENGR 2405 -- Electrical Circuits I}\\[2pt]
    {\large\bfseries Laboratory \#<N>: <Lab Title>}\\[10pt]
    \begin{tabular}{rl}
        Name:    & Abdon Morales \\
        UT EID:  & am226923 \\
        Date:    & \today \\
    \end{tabular}
\end{center}
\vspace{0.5em}
\hrule
\vspace{1.5em}

% ----------------------------- description -----------------------------------
\section*{Description}
<Brief description of the lab, taken from the lab assignment and put in your own
words. State what circuit is built, what is measured, and what the lab is meant
to demonstrate.>

%==============================================================================
\section*{Procedure}
%==============================================================================
\begin{enumerate}[leftmargin=*]
    \item <First step.>
    \item <Second step.>
    \item <Third step.>
\end{enumerate}

%==============================================================================
\section*{Calculations}
%==============================================================================
% Delete this section if the lab has no hand calculations.
<Set up the governing equations before any numbers go in.>
\[
    V = IR
\]

% Pull the program straight out of the .sce file (uncomment when the file exists):
% \lstinputlisting[style=scilab, caption={SciLab program \texttt{lab.sce}.}]{lab.sce}

% Or paste it inline:
\begin{lstlisting}[style=scilab, caption={SciLab program.}]
// paste program here
R = [10 20 30];
Rseries = sum(R)
\end{lstlisting}

\begin{lstlisting}[style=console, caption={SciLab runtime output.}]
 Rseries =
    60.
\end{lstlisting}

%==============================================================================
\section*{Simulations}
%==============================================================================
\begin{figure}[H]
    \centering
    % \includegraphics[width=0.85\textwidth]{schematic.png}
    \fbox{\parbox[c][3.5cm][c]{0.85\textwidth}{\centering <LTspice schematic image>}}
    \caption{LTspice schematic.}
\end{figure}

\begin{lstlisting}[style=netlist, caption={LTspice netlist.}]
* paste netlist here
V1 N001 0 SINE(0 5 1k)
R1 N001 0 1k
.tran 5m
.end
\end{lstlisting}

%==============================================================================
\section*{Measurements}
%==============================================================================
\begin{table}[H]
    \centering
    \caption{Calculated and measured values.}
    \begin{tabular}{lccc}
        \toprule
        Quantity & Calculated & Measured & \% Error \\
        \midrule
        $V_1$ & <\ > & <\ > & <\ > \\
        $I_1$ & <\ > & <\ > & <\ > \\
        \bottomrule
    \end{tabular}
\end{table}

\begin{figure}[H]
    \centering
    % \includegraphics[width=0.8\textwidth]{scope_shot.png}
    \fbox{\parbox[c][3.5cm][c]{0.8\textwidth}{\centering <Oscilloscope capture>}}
    \caption{Oscilloscope capture.}
\end{figure}

%==============================================================================
\section*{Observations \& Discussion}
%==============================================================================
<Connected prose. Compare measured against calculated, explain the size and
direction of any discrepancy, flag anything anomalous in the data, and answer
the questions asked in the lab handout.>

%==============================================================================
\end{document}
